\begin{center}
\begin{adjustbox}{max width=.82\textwidth}
\begin{threeparttable}[!h]
\caption{Women's Migration Patterns Based on Characteristics of Origin and Destination Districts}
\label{tab:flows_orig_dest}
\begin{tabular}{lccc}
\toprule
&\multicolumn{1}{c}{(1)}&\multicolumn{1}{c}{(2)}&\multicolumn{1}{c}{(3)} \\
{ \scshape A1. Districts characteristics}&\makecell[c]{\scshape  Rural}& \makecell[c]{\scshape  Urban}& \makecell[c]{\scshape  Total} \\
\midrule
Number of districts &       164 &       104 &       268 \\
Migration rate &      0.16 &      0.21 &      0.18 \\
Mean district FLFP rate &      0.58 &      0.47 &      0.54 \\
SD district FLFP rate &      0.14 &      0.11 &      0.14 \\
\midrule & \multicolumn{3}{c}{\scshape Birth recency urbanicity} \\  \cmidrule(lr){2-4}  \scshape A2. Share currently residing in &\scshape Rural & \scshape Urban &\scshape Total \\ \midrule
Urban districts &      0.56 &      0.70 &      0.62 \\
Rural districts &      0.44 &      0.30 &      0.38 \\
\midrule\\ \scshape B1.  Districts characteristics & \scshape Low-FLFP & \scshape High-FLFP & \scshape Total \\ \midrule 
Number of districts &       109 &       159 &       268 \\
Migration rate &      0.18 &      0.18 &      0.18 \\
Mean district FLFP rate &      0.40 &      0.63 &      0.54 \\
SD district FLFP rate &      0.06 &      0.09 &      0.14 \\
\midrule & \multicolumn{3}{c}{\scshape Birth recency FLFP} \\ \cmidrule(lr){2-4} \scshape B2. Share currently residing in &\scshape Low-FLFP &\scshape High-FLFP &\scshape Total \\ \midrule
High-FLFP districts &      0.22 &      0.48 &      0.38 \\
Low-FLFP districts &      0.78 &      0.52 &      0.62 \\
\bottomrule
\bottomrule
\end{tabular}
\begin{tablenotes}
\item \footnotesize \textit{Notes:} The table summarizes migration flows from the pooled 1985, 1995, and 2005 Intercensal Surveys based on the characteristics of both origin and destination districts. I used data from the 2010 Indonesian Census to classify districts by urbanicity and female labor force participation (FLFP) rates. I classified districts by urbanicity and FLFP rates using data from the 2010 Indonesian census. In panels A1 and A2, I classify a district as urban if more than 43.3% of its population resides in areas designated as urban by the Indonesian Central Bureau of Statistics (BPS). This threshold was selected to ensure that the proportion of the population living in urban districts aligns with the national urban population share reported by BPS. In panels B1 and B2, I split districts at the median FLFP rate. 
\end{tablenotes}
\end{threeparttable}
\end{adjustbox}
\end{center}
